% ============================================================
% Introducción (incluye problema, objetivos, justificación)
% ============================================================
\chapter{Introducción}
\label{cap:introduccion}

% ============================================================
% >>> NOTA-GUÍA DE AUTORÍA (borrar antes de entregar) <<<
% TODO ESTE CAPÍTULO LO ESCRIBÍS VOS, con tus palabras.
%   - El planteamiento, las preguntas y los objetivos de abajo son un BORRADOR
%     basado en tus ideas: reescribilos en tu voz y verificá que digan lo que
%     realmente vas a hacer.
%   - No requiere citas obligatorias, pero si afirmás datos del contexto
%     (p. ej. uso de SO, estadísticas) citá la fuente en APA: \parencite{clave}.
% ============================================================

Presentación general del tema, contexto en el que se inscribe el trabajo y
relevancia del estudio para la carrera de \carreraTesis. La introducción ofrece
al lector una visión panorámica del problema abordado y de la estructura del
documento.

\section{Planteamiento del problema}

Todo sistema operativo de propósito general moderno emplea planificación
\emph{apropiativa}: el núcleo utiliza las interrupciones de temporizador
(hardware) para desalojar al proceso en ejecución al expirar su \emph{quantum} y
conmutar a otro listo. Sin embargo, dentro de este paradigma coexisten distintas
\emph{políticas} de planificación —Round Robin (RR), colas multinivel con
retroalimentación (\emph{Multilevel Feedback Queue}, MLFQ) y planificación por
prioridades— que representan compromisos de diseño diferentes. RR ofrece un
reparto equitativo y un tiempo de respuesta acotado y predecible, con una
implementación simple y de bajo costo, lo que lo hace atractivo para sistemas de
recursos limitados o con requisitos de previsibilidad. MLFQ y las políticas por
prioridades, en cambio, favorecen a las tareas interactivas y mejoran el tiempo
de respuesta promedio en cargas mixtas, a costa de una mayor complejidad y del
riesgo de inanición (\emph{starvation}). El problema consiste en determinar,
mediante medición sobre un núcleo real, en qué medida cada política se ajusta
mejor a distintos contextos de uso y condiciones de carga.

\subsection{Formulación del problema}
¿En qué medida las distintas políticas de planificación apropiativa (Round
Robin, colas multinivel con retroalimentación y prioridades), implementadas y
medidas sobre un núcleo real, difieren en la distribución del tiempo de CPU, el
tiempo de respuesta, la productividad, la inanición (\emph{starvation}) y el
costo de cambio de contexto, y qué contexto de uso favorece a cada una?

\subsection{Preguntas específicas}
\begin{itemize}
  \item ¿Cómo incide el tamaño del \emph{quantum} en el tiempo de respuesta, el
    tiempo de espera promedio y el costo de cambio de contexto?
  \item ¿Qué diferencias de productividad (\emph{throughput}), latencia y
    equidad presentan Round Robin y MLFQ bajo cargas con mezcla de tareas
    interactivas y de cómputo?
  \item ¿Bajo qué condiciones aparece la inanición en MLFQ y en qué medida la
    corrige un mecanismo de promoción periódica (\emph{priority boost})?
  \item ¿Qué fracción del tiempo de CPU consume el propio mecanismo de
    planificación (sobrecarga de cambio de contexto) en cada política y cómo
    escala con el número de procesos?
  \item ¿Qué política resulta más adecuada para contextos de previsibilidad o
    recursos limitados frente a contextos de multitarea interactiva?
\end{itemize}

\section{Objetivos}

\subsection{Objetivo general}
Analizar y evaluar comparativamente el rendimiento de distintas políticas de
planificación apropiativa de procesos —Round Robin, colas multinivel con
retroalimentación y prioridades— en el núcleo de un sistema operativo, a partir
de su implementación y medición sobre un prototipo real escrito en lenguaje C,
con el fin de determinar qué política se ajusta mejor a cada contexto de uso.

\subsection{Objetivos específicos}
\begin{itemize}
  \item Caracterizar los fundamentos y el comportamiento de las principales
    políticas de planificación apropiativa de procesos (Round Robin, MLFQ y
    prioridades) en el núcleo de un sistema operativo.
  \item Definir las métricas e indicadores de rendimiento (tiempo de respuesta,
    tiempo de espera, tiempo de retorno, productividad, inanición y
    sobrecarga de cambio de contexto) que servirán de base para la
    comparación.
  \item Implementar dichas políticas sobre el núcleo, en lenguaje C, utilizando
    su mecanismo de interrupciones de temporizador y de cambio de contexto como
    instrumento de medición.
  \item Medir y comparar el rendimiento de cada política bajo distintas
    condiciones de carga y tamaños de \emph{quantum}, incluyendo el costo del
    cambio de contexto.
  \item Interpretar los resultados para determinar qué política conviene en
    contextos de previsibilidad y recursos limitados frente a contextos de
    multitarea interactiva.
\end{itemize}

\section{Justificación}

Razones que fundamentan la importancia y pertinencia del estudio (teórica,
práctica, social, tecnológica).

\section{Alcance y limitaciones}

Delimitación del estudio (temporal, espacial, temática) y restricciones
encontradas.
